\documentclass[11pt,letterpaper]{article}
\usepackage[utf8]{inputenc}
\usepackage[T1]{fontenc}
\usepackage{lmodern}
\usepackage[margin=1in]{geometry}
\usepackage{amsmath,amssymb,amsthm,mathtools}
\usepackage{booktabs,array,longtable}
\usepackage{microtype}
\usepackage[unicode,hidelinks]{hyperref}
\setlength{\emergencystretch}{3em}
\newtheorem{proposition}{Proposition}[section]
\newtheorem{definition}[proposition]{Definition}
\newtheorem{assumption}[proposition]{Assumption}
\newtheorem{remark}[proposition]{Remark}
\newcommand{\R}{\mathbb R}
\newcommand{\E}{\mathbb E}
\newcommand{\Pp}{\mathbb P}
\newcommand{\F}{\mathcal F}
\newcommand{\dd}{\,\mathrm d}
\newcommand{\ind}{\mathbf 1}
\DeclareMathOperator{\diag}{diag}
\title{Commodities and Money, \\ working research notes for\\a mathematical reconstruction of \\ \emph{Capital},\\Volume I, Part I}
\author{Felix Alfonso Camacho Moya}
\date{October 13, 2026}
\begin{document}
\maketitle
\begin{abstract}
We begin a chapter-by-chapter mathematical research program on Marx's
\emph{Capital}. This installment concerns Volume I, Part I, Chapters 1--3.
It separates physical quantities, labor-value magnitudes, value expressions,
and monetary realization. A finite production network supplies one explicit
closure for socially necessary labor coefficients; a graph formalizes
consistent quantitative expressions; a commodity-money standard converts
labor magnitudes into monetary magnitudes; and a constrained stochastic
exchange process separates a quoted price from an accomplished sale.
We prove conditional results on production, exchange consistency,
conservation, and realization delay, and derive stochastic comparative
statics. These are results of the stated models, not proofs of Marx's
theory or of historically inevitable crises.
\end{abstract}
\tableofcontents

\section{Scope, sources, and method}
The primary textual sequence is Chapter 1 (commodities), Chapter 2
(exchange), and Chapter 3 (money) \cite{marx1,marx2,marx3}. Chapter and
section references are used because pagination varies between editions.
The library copy \texttt{Capital-Volume-I.pdf} identifies the Moore--Aveling
translation edited by Engels and supplies a local reference edition.

The purpose is a reconstruction with explicit commitments. A \emph{textual
category} is something to interpret; an \emph{assumption} selects an
operational representation; a \emph{proposition} follows from that
representation. We do not infer that a category has been explained merely
because a variable has been assigned to it. In particular, neither the
historical institution of money nor the reduction of heterogeneous labor
is derived by an algebraic change of units.

The existing \emph{Stochastic Marxian Macroeconomics} draft uses
$\lambda=\ell(I-A)^{-1}$ and a monetary labor-time conversion in its
production network. Here we examine the assumptions behind that earlier
notation before introducing surplus-value, wages, accumulation, or
financial claims. Those categories require later installments.

\begin{center}
\small
\begin{tabular}{p{0.16\linewidth}p{0.34\linewidth}p{0.40\linewidth}}
\toprule
Textual location & Category to retain & Mathematical representation here\\
\midrule
Ch. 1, §§1--2 & Useful goods; concrete and abstract labor & Typed quantities;
labor-reduction coefficients; a conditional production closure\\
Ch. 1, §3 & Relative and equivalent forms & Oriented value expressions;
consistent edge ratios; a designated common equivalent\\
Ch. 1, §4 & Fetishism & A many-to-one map from social data to observable
magnitudes; an explicit limit of the representation\\
Ch. 2 & Owners and exchange & Ownership inventories and admissible
bilateral transfers\\
Ch. 3, §§1--2 & Measure and circulation & Commodity-money standard;
quoted prices; completed sales and monetary balances\\
Ch. 3, §3 & Hoards, payment, world money & A circulation extension and
an agenda for contracts and international settlement\\
\bottomrule
\end{tabular}
\end{center}

\section{Chapter 1: quantities and the two aspects of labor}
\subsection{Goods are typed quantities, not interchangeable numbers}
Fix $n$ reproducible commodity types and $H$ owners. A type includes a
quality specification, location and delivery date when relevant, and a
physical unit. A physical bundle belongs to
\[
 X=\prod_{i=1}^n\R_+\,e_i,
\]
where $e_i$ labels the unit of type $i$. Owner $a$ holds $x^a\in X$.
The addition of coordinates records inventories; it does not assert that
a meter of cloth is physically comparable to a kilogram of grain.

Let $U_{ai}(q,z)\in\{0,1\}$ indicate whether $q$ units of type $i$ can
serve a relevant use for owner $a$ in social state $z$. This predicate is
not a cardinal utility function. It permits usefulness to depend on
quality, infrastructure, and social circumstances. Useful objects need
not have labor value, and possession of an object does not guarantee a
buyer. The model below covers reproducible products of labor offered for
exchange. It does not force land or every priced object into that domain.

\subsection{Concrete activity and a social reduction operator}
Let $k=1,\ldots,K$ label labor activities. Suppose $L_{kj}$ is the time of
activity $k$ required per unit of output $j$ under the selected normal
production conditions. Specify positive reduction weights
$\alpha_k(z)$, measured in simple-labor hours per hour of activity $k$.
The resulting direct abstract-labor coefficient is
\begin{equation}\label{eq:reduction}
 \ell_j(z)=\sum_{k=1}^K\alpha_k(z)L_{kj}(z).
\end{equation}
Concrete activities remain distinct in $L$; $\ell$ records their reduction
to one accounting magnitude. Linearity and the weights are assumptions.
They do not establish the social process by which complex labor is
reduced to simple labor, nor explain why privately performed labor
acquires social validity.

For a firm $f$, let $h_{fj}$ be its private hours per unit. We do not set
the social coefficient equal to $h_{fj}$. A producer's inefficiency
cannot increase the social magnitude simply by consuming more time.
One possible empirical rule is an output-weighted mean over a specified
set of normal producers. Other rules select regulating techniques or
normal capacity. Choosing between them is a substantive research
question; a mean is not supplied by the phrase ``socially necessary.''

\subsection{A simultaneous production closure}
\begin{assumption}[Normal reproducible technology]\label{ass:technology}
Fix physical units and one accounting period. Let
$A=(a_{ij})\in\R_+^{n\times n}$, where $a_{ij}$ is units of input $i$
consumed per unit of output $j$, and let
$\ell\in\R_{++}^{1\times n}$ be direct simple-labor hours per output
unit. Assume the spectral radius $\rho(A)<1$.
We use circulating inputs, fixed coefficients, no joint production,
and no endogenous technique choice. The coefficients describe socially
normal production, not an arbitrary individual workshop.
\end{assumption}

\begin{definition}[Reproduction labor coefficients]\label{def:value}
A row vector $v$ of labor hours per commodity unit satisfies
\begin{equation}\label{eq:value}
 v=vA+\ell.
\end{equation}
The labor magnitude assigned to bundle $q$ is $V(q)=vq$.
\end{definition}
Equation \eqref{eq:value} is an input-output reconstruction, not an
equation written by Marx in Part I. Its recursive treatment of consumed
inputs anticipates analysis of production. It identifies $v$ with
socially necessary labor requirements only within the chosen closure.
An alternative temporal closure values inputs at their production or
acquisition dates; these alternatives need not agree during technical
change. We retain the simultaneous closure as a transparent baseline.

\begin{proposition}[Existence, uniqueness, and positivity]\label{prop:values}
Under Assumption \ref{ass:technology}, the unique solution is
\begin{equation}\label{eq:inverse}
 v=\ell R,\qquad R=(I-A)^{-1}=\sum_{r=0}^{\infty}A^r,
\end{equation}
and $v_i>0$ for every $i$.
\end{proposition}
\begin{proof}
The spectral-radius assumption makes the Neumann series converge to the
inverse of $I-A$. Every summand is nonnegative and $R\ge I$ entrywise.
Thus $\ell R\ge\ell>0$. Multiplying \eqref{eq:value} on the right by
$R$ proves existence and uniqueness. The term $\ell A^r$ counts labor
through input chains of length $r$.
\end{proof}

\begin{proposition}[Net-product identity]\label{prop:net}
For a nonnegative gross output $x$ with net product $y=(I-A)x\ge0$,
\begin{equation}\label{eq:net}
 vy=\ell x.
\end{equation}
\end{proposition}
\begin{proof}
Multiply $v(I-A)=\ell$ by $x$.
\end{proof}
The identity concerns the \emph{net} product. Gross output has labor
magnitude $vx=vAx+\ell x$; confusing it with newly performed labor
double-counts intermediate inputs. The identity is conditional on the
production closure and does not prove that market revenue equals it.

\begin{proposition}[Comparative statics]
For infinitesimal changes of admissible coefficients,
\begin{equation}\label{eq:sensitivity}
 \delta v=(\delta\ell+v\,\delta A)R.
\end{equation}
If $A'\le A$ and $\ell'\le\ell$ entrywise, with both systems satisfying
the assumptions, then $v'\le v$ entrywise.
\end{proposition}
\begin{proof}
Differentiate $v(I-A)=\ell$. For the finite comparison, induction gives
$(A')^r\le A^r$ for every $r$. Apply the nonnegative Neumann series.
\end{proof}
The proposition specifies one sense in which greater productivity
reduces labor requirements. Technique substitution, quality changes,
joint products, and changing labor-reduction weights require new
arguments. More physical output need not imply more labor magnitude.

\section{Chapter 1: value expressions and their quantitative geometry}
\subsection{The simple expression is oriented}
The expression
\[
 q_i\text{ units of commodity }i
 \quad\longmapsto\quad q_j\text{ units of commodity }j
\]
places $i$ in the relative role and $j$ in the equivalent role. Under an
ideal expression at the labor coefficients above,
\begin{equation}\label{eq:ratio}
 q_iv_i=q_jv_j,\qquad r_{ij}=q_j/q_i=v_i/v_j.
\end{equation}
The scalar equality is symmetric; the roles in the expression are not.
Store the role assignment as an oriented edge. Algebraic reciprocity
$r_{ji}=r_{ij}^{-1}$ does not itself imply reciprocal direct saleability.
Real market quotations are introduced separately below.

An expanded expression for one commodity is a family of outgoing
edges to several equivalents. A general expression arranges many
commodities against a common equivalent $g$. These graph structures
describe quantitative forms; they do not establish a historical sequence
of barter stages.

\subsection{When do ratios have a common scalar representation?}
Consider a connected undirected graph whose edges have strictly positive
ratios in each orientation, with $r_{ji}=r_{ij}^{-1}$. The graph here
records comparability, not permission to execute a trade.
\begin{proposition}[Cycle consistency]\label{prop:cycles}
There is a positive vector $w$, unique up to a common positive factor,
such that $r_{ij}=w_i/w_j$ on every edge if and only if the product of
ratios around every cycle is $1$.
\end{proposition}
\begin{proof}
If such $w$ exists, multiplication telescopes around any cycle.
Conversely fix a root $g$ and $w_g=1$. Define $w_i$ as the product of
edge ratios along a path from $i$ to $g$. Any two paths give the same
product because their concatenation is a closed walk, whose factors
cancel by the cycle condition. For the edge $i\to j$, this gives
$w_i=r_{ij}w_j$. Connectivity fixes all coordinates relative to $g$.
\end{proof}

This theorem does \emph{not} identify $w$ as labor values. Any consistent
positive price vector gives such ratios. The labor interpretation enters
through \eqref{eq:value}, independent of graph consistency.
For disconnected graphs, there is one undetermined scale per component.

Taking logarithms gives
\[
 b_{ij}=\log r_{ij}=\phi_i-\phi_j,\qquad \phi_i=\log w_i.
\]
Thus edge data must be an exact potential difference. With noisy data,
one can estimate $\phi$ by weighted least squares over the edges,
fixing $\phi_g=0$. Nonzero cycle sums are quantitative inconsistencies
to investigate. They need not be executable arbitrage: accessibility,
timing, inventory, bid--ask spreads, and transaction costs matter.

\subsection{A common equivalent needs an institution}
\begin{assumption}[Money institution]\label{ass:money}
One commodity $g$ is socially accepted in settlement in a specified
domain. It is divisible at the scale of this model, has positive labor
coefficient $v_g$, and a fixed monetary denomination $\nu>0$, measured
in currency units per physical unit of $g$.
\end{assumption}
The choice of $g$ is an institutional assumption. A graph always permits
a mathematical root; this is insufficient to explain the selection,
acceptance, or durability of money. Commodity money is the baseline
appropriate to the text's gold abstraction. Fiat money and bank deposits
require a distinct institutional model.

\section{Chapter 1: a limited representation of fetishism}
Let $z$ contain techniques, labor activities, ownership, contractual
rules, and patterns of social acceptance. Observable magnitudes may be
represented by a map
\[
 \Psi:z\longmapsto\bigl(v(z),p(z)\bigr).
\]
This map is usually many-to-one. For example, with one sector and labor
already reduced to simple hours,
\[
 (A,\ell)=(0,1)\quad\text{and}\quad(A,\ell)=(1/2,1/2)
\]
both give $v=1$. Ownership changes can also leave $A$, $\ell$, and a
chosen price rule unchanged. Magnitudes alone therefore do not identify
the relations represented by the full state.

This non-identification is a useful mathematical warning against
treating a price as an intrinsic physical attribute. It is not a full
theory of fetishism: a statistical projection cannot by itself explain
the objective social power of commodity relations or the appearance
those relations assume for participants. An adequate extension must
retain ownership and institutions in the state rather than discard them
after computing $v$.

\section{Chapter 3: labor magnitudes and monetary expression}
\subsection{Commodity-money standard and value prices}
Under Assumption \ref{ass:money}, define
\begin{equation}\label{eq:melt}
 m=\frac{\nu}{v_g},\qquad p_i^*=mv_i=\nu\frac{v_i}{v_g}.
\end{equation}
Here $m$ is currency units per simple-labor hour and $p_i^*$ is currency
units per commodity unit. This is a \emph{commodity-standard monetary
expression of labor time}; it need not coincide with an aggregate MELT
defined from monetary net output and living labor in a different model.
If $g$ is included among the $n$ commodities, \eqref{eq:melt} gives
$p_g^*=\nu$. If $g$ is external, its production coefficient must be
supplied by a separate system; it cannot be inferred from $A$.

Fixing a currency unit measures prices; specifying $v_g$ relates that
standard to labor magnitudes. A denomination change $\nu\mapsto c\nu$
multiplies all $p_i^*$ and $m$ by $c$ without changing physical quantities
or labor coefficients. If all $v_i$, including $v_g$, fall by the same
factor, the value prices $p_i^*$ remain unchanged even though labor
requirements have fallen.

\subsection{Quoted prices and realized receipts}
Introduce a positive quoted-price vector $p$ independently. For the
reproducible labor products in our domain, parameterize
\begin{equation}\label{eq:wedge}
 p_i=p_i^*\exp(\xi_i).
\end{equation}
This is a decomposition, not a theory of $\xi_i$. If $g$ is included,
set $\xi_g=0$ to preserve the denomination. Scarcity, bargaining,
institutional price setting, and market access can later determine the
wedge. We do not impose $\sum_i p_iq_i=m\sum_i v_iq_i$ for arbitrary
traded bundles; such an aggregate equality needs an additional closure.

If units actually sold during an interval are recorded by finite-variation
processes $Q_i^{\rm sold}$, then gross receipts are
\begin{equation}\label{eq:receipts}
 R_T^{\rm sales}=\sum_i\int_{(0,T]}p_{i,t-}\dd Q_{i,t}^{\rm sold}.
\end{equation}
A quote or an inventory marked at that quote is not this cash flow.
Products can remain unsold when every displayed $p_i$ is positive.

\section{Chapter 2 and Chapter 3: a stochastic exchange model}
\subsection{State and admissible events}
Work on a filtered probability space
$(\Omega,\F,(\F_t)_{t\ge0},\Pp)$ satisfying the usual conditions.
Owner $a$ holds non-money inventory $x^a_t\in\R_+^{n_c}$ and money
balance $b^a_t\ge0$, where $n_c$ counts non-money types.
Production is paused during the exchange interval. Balances are claims
to the accepted settlement asset; for the commodity-money baseline,
$b^a/\nu$ is the physical stock of $g$ owned by $a$.

A sale event $e=(a,c,i,q)$ means seller $a$ transfers $q>0$ units of
commodity $i$ to buyer $c\ne a$ at the pre-event quote $p_i$. Its updates
are
\begin{equation}\label{eq:trade}
\begin{aligned}
 x^{a\prime}&=x^a-qe_i,& x^{c\prime}&=x^c+qe_i,\\
 b^{a\prime}&=b^a+p_iq,& b^{c\prime}&=b^c-p_iq.
\end{aligned}
\end{equation}
It is admissible only if $x_i^a\ge q$, $b^c\ge p_iq$, the buyer can
use or pass on the good, and both parties accept the transfer. Those
conditions represent ownership and contract restrictions. They do not
reduce exchange to an automatic physical motion of objects.

For a finite list of possible lot sizes, let $N^e_t$ count sales of type
$e$ with predictable intensity $\kappa_e(t,Z_{t-})$, where $Z$ includes
inventories, balances, quotes, and acceptance variables. Set the intensity
to zero outside the admissible set. Suppose the total intensity is
bounded on each finite time horizon and events do not occur
simultaneously. A construction by thinning marked Poisson events then
produces a nonexplosive process. Conditional on an exogenous quote path,
the exchange jump operator is
\begin{equation}\label{eq:generator}
 (\mathcal J_t f)(z)=\sum_e\kappa_e(t,z)
       \bigl[f(T_e(t,z))-f(z)\bigr].
\end{equation}
If quotes themselves have Markov dynamics, their generator must be
added for a generator of the full state. Formula \eqref{eq:generator}
alone describes the exchange part.

\begin{proposition}[Nonnegativity and conservation]\label{prop:conservation}
The sale process preserves nonnegative holdings and, pathwise,
\begin{equation}\label{eq:conservation}
 \sum_a x^a_{i,t}=\sum_a x^a_{i,0},\qquad
 \sum_a b^a_t=\sum_a b^a_0.
\end{equation}
For labor coefficients frozen during the interval, total inventory labor
magnitude is also conserved.
\end{proposition}
\begin{proof}
Admissibility prevents a seller or buyer from giving more than it holds.
Each update subtracts exactly what it adds to another owner's account.
Sum the updates over the finite number of events on each finite horizon.
Multiply the physical conservation identity by the fixed coefficients
to obtain conservation of inventory labor magnitude.
\end{proof}
Exchange can redistribute money and goods without creating total labor
magnitude in this model. It need not preserve an individual owner's
labor-equivalent holdings when $p_i\ne p_i^*$. Nor does the proposition
exclude changing monetary wealth when quotes move: revaluation is
different from production. With changing $v_t$, inventory conservation
does not imply constant $v_tx_t$.

\subsection{The two acts of simple circulation}
For an owner selling type $i$ to buy type $j$, let $\tau_s$ be the first
successful sale and $\tau_b\ge\tau_s$ the subsequent intended purchase.
The circuit is
\[
 C_i\xrightarrow{\tau_s}M\xrightarrow{\tau_b}C_j.
\]
The purchase may never occur. Between the acts, money can remain idle.
Its purchasing power over $j$ changes with $p_{j,t}$. The two stopping
times therefore retain a separation lost in instantaneous barter.
This is simple commodity circulation; it does not yet describe the
$M$--$C$--$M'$ circuit of capital.

\begin{proposition}[Sale delay under a specified hazard]\label{prop:delay}
Suppose a single lot has no competing sale event and has deterministic,
locally integrable first-sale hazard $h(t)\ge0$. Then
\begin{equation}\label{eq:survival}
 \Pp(\tau_s>T)=\exp\left(-\int_0^T h(u)\dd u\right),
 \qquad
 \E[\tau_s]=\int_0^\infty
       \exp\left(-\int_0^t h(u)\dd u\right)\dd t,
\end{equation}
where the expectation may be infinite. For constant $h>0$,
$\E[\tau_s]=1/h$.
\end{proposition}
\begin{proof}
The survival function solves $S'(t)=-h(t)S(t)$ almost everywhere with
$S(0)=1$. The expectation formula follows by integrating the tail of
a nonnegative random variable.
\end{proof}
For $h\equiv0$, a positive quote gives no realization. Even when
$\int_0^\infty h(u)\dd u=\infty$ guarantees eventual sale almost surely,
mean delay can be infinite: $h(t)=1/(1+t)$ gives $S(t)=1/(1+t)$.
Thus almost-sure realization and timely realization are different.

For a random hazard independent of the exponential trigger, conditional
on the hazard path the same survival formula holds, and the unconditional
survival is its expectation. That Cox-process formula must not be
applied automatically to an endogenous hazard that feeds back on the
sale itself. A general predictable intensity gives the compensator
identity
\begin{equation}\label{eq:compensator}
 N^e_t-\int_0^t\kappa_e(u,Z_{u-})\dd u
 \quad\text{is a local martingale}.
\end{equation}
With suitable integrability it is a martingale. Expected receipts can
then be calculated from compensators, for example
\[
 \E[R_T^{\rm sales}]
 =\E\left[\sum_e\int_0^T
       p_{i(e),t-}q_e\kappa_e(t,Z_{t-})\dd t\right],
\]
when the nonnegative sums and integrals are well defined (or integrable
for a finite-valued expectation).

\subsection{Circulation, hoarding, and a finite accounting identity}
Divide accepted money into active balances $B_t^{\rm act}$ and hoards
$H_t^{\rm hoard}$, so $B=B_t^{\rm act}+H_t^{\rm hoard}$ in this closed
exchange system. Hoarding changes availability and potentially the
intensities; it does not destroy the accepted money stock.

Let $F_{[0,T]}$ be the sum of monetary amounts in completed cash sales.
For $T>0$ and mean active balance
$\overline B^{\rm act}=T^{-1}\int_0^T B_t^{\rm act}\dd t>0$, define
the turnover count $V_T=F_{[0,T]}/\overline B^{\rm act}$. Then
\begin{equation}\label{eq:turnover}
 F_{[0,T]}=\overline B^{\rm act}V_T.
\end{equation}
This is an accounting definition. It supplies no causal rule that a
money-stock change must cause a proportional price change. Moreover,
$F$ is realized sales, not the sum of prices of all goods offered.
Credit sales, debt settlement, netting, and cross-border payments require
an expanded ledger before they can be included.

\section{Stochastic production coefficients and value prices}
\subsection{A precise local domain}
Let $W$ be a $d$-dimensional Brownian motion under $\Pp$. On an open
domain with $I-A$ invertible, suppose
\begin{equation}\label{eq:coefficients}
\begin{aligned}
 \dd A_t&=B_t\dd t+\sum_{k=1}^d C_t^k\dd W_t^k,\\
 \dd\ell_t&=a_t\dd t+\sum_{k=1}^d c_t^k\dd W_t^k.
\end{aligned}
\end{equation}
Here $B,C^k$ are matrices and $a,c^k$ are row vectors. Assume adapted
coefficients with the usual local integrability for these It\^o
processes. For the economic interpretation, require $A_t\ge0$,
$\ell_t>0$, and $\rho(A_t)<1$.

An unrestricted Brownian coefficient model need not preserve those
inequalities. The following calculation is local, up to the first exit
from the admissible domain, with localization bounding $R_t$ and the
coefficients. A global economic model needs an invariant-domain
construction, reflection, or suitable transformations. A bound on
$\rho(A)$ alone is not an entrywise bound on $A$ or a uniform bound on
the inverse for all nonnormal matrices.

\subsection{It\^o formula for network labor coefficients}
Set $R_t=(I-A_t)^{-1}$ and $v_t=\ell_tR_t$. Differentiating the inverse
and using matrix It\^o products gives
\begin{equation}\label{eq:dR}
\begin{aligned}
 \dd R_t={}&\left(R_tB_tR_t+
      \sum_kR_tC_t^kR_tC_t^kR_t\right)\dd t
      +\sum_kR_tC_t^kR_t\dd W_t^k.
\end{aligned}
\end{equation}
Consequently
\begin{equation}\label{eq:dv}
\begin{aligned}
 \dd v_t={}&\left[
     a_tR_t+v_tB_tR_t
     +\sum_k(v_tC_t^k+c_t^k)R_tC_t^kR_t
             \right]\dd t\\
 &+\sum_k(c_t^k+v_tC_t^k)R_t\dd W_t^k.
\end{aligned}
\end{equation}
To verify the matrix order, start from $(I-A_t)R_t=I$. The diffusion
coefficient of $R$ is $RC^kR$. Its quadratic cross-product with $A$
then yields the drift correction $RC^kRC^kR$. In
$\dd(\ell R)=\dd\ell\,R+\ell\,\dd R+\dd\ell\,\dd R$, the last
term is $\sum_kc^kRC^kR\dd t$, which supplies the labor--technology
covariance correction in \eqref{eq:dv}.

These are exact local identities for the simultaneous closure.
They do not supply the coefficients in \eqref{eq:coefficients}, nor
derive a law of technical change from exchange. A stochastic law of
motion must follow an economic specification, rather than stand in for
one.

\subsection{Relative prices when money also has a changing labor value}
Suppose positive coefficients $v_i,v_g$ have local dynamics
\begin{equation}\label{eq:relativevalue}
 \frac{\dd v_i}{v_i}=\beta_i\dd t+\sigma_i\cdot\dd W_t,
 \qquad
 \frac{\dd v_g}{v_g}=\beta_g\dd t+\sigma_g\cdot\dd W_t,
\end{equation}
with all coefficients evaluated at $t$. For fixed $\nu$, the ideal
commodity-money price $p_i^*=\nu v_i/v_g$ satisfies
\begin{equation}\label{eq:priceSDE}
 \frac{\dd p_i^*}{p_i^*}
 =\bigl[\beta_i-\beta_g+\|\sigma_g\|^2
                    -\sigma_i\cdot\sigma_g\bigr]\dd t
       +(\sigma_i-\sigma_g)\cdot\dd W_t.
\end{equation}
Equivalently,
\begin{equation}\label{eq:logprice}
 \dd\log p_i^*
 =\left[\beta_i-\beta_g-\tfrac12
                  (\|\sigma_i\|^2-\|\sigma_g\|^2)\right]\dd t
       +(\sigma_i-\sigma_g)\cdot\dd W_t.
\end{equation}
The inverse-money correction and the covariance term cannot be dropped.
Identical proportional dynamics for $v_i$ and $v_g$ leave their ratio
constant. If $g$ belongs to the network, its dynamics are the corresponding
coordinate of \eqref{eq:dv}; they cannot be chosen independently.

For continuous It\^o $\xi_i$ in \eqref{eq:wedge},
$\dd\log p_i=\dd\log p_i^*+\dd\xi_i$. Converting to $\dd p_i/p_i$
requires the quadratic variation of this entire sum, including cross
covariation. This is one route to a model of price fluctuations anchored
in production, but neither $p_i^*$ nor $p_i$ is assumed to be a
risk-neutral asset price.

\section{Worked examples and diagnostic counterexamples}
\subsection{Two reproducible goods and an external money commodity}
Take
\[
 A=\begin{pmatrix}1/10&1/5\\1/5&1/10\end{pmatrix},
 \qquad \ell=(2/5,4/5).
\]
The eigenvalues are $3/10$ and $-1/10$, so the assumptions hold. Direct
calculation gives
\[
 R=\frac1{77}\begin{pmatrix}90&20\\20&90\end{pmatrix},
 \qquad v=(52/77,80/77).
\]
For net output $y=(10,5)^\top$,
\[
 x=Ry=(1000/77,650/77)^\top,
 \qquad vy=\ell x=920/77\text{ labor hours}.
\]
Specify an external money commodity with $v_g=1/2$ hour per unit and
$\nu=10$ currency units per unit. Then $m=20$ and
\[
 (p_1^*,p_2^*)=(1040/77,1600/77)
       \approx(13.5065,20.7792).
\]
The ideal ratio is $r_{12}=13/20$. Halving $\ell$ with fixed $A$ and
fixed $v_g$ halves $v$ and these prices. Halving both $\ell$ and $v_g$
instead leaves these prices unchanged. The money commodity's production
specification is additional data in this example.

\subsection{Cycle inconsistency does not identify its social cause}
Three observed oriented ratios $r_{12}=2$, $r_{23}=3$, and $r_{31}=1/5$
have cycle product $6/5$. No scalar vector reproduces all three
simultaneously. This diagnoses failure of the scalar representation;
it does not tell us whether the observations differ in delivery date,
quality, access, or bargaining conditions.

\subsection{Money available in aggregate need not be available to a buyer}
Owners $A,B,C$ begin with money balances $(6,0,0)$. Suppose $B$ owns
one unit of good 1 and $C$ one unit of good 2, each priced at 6. First
$A$ buys good 1 from $B$, producing balances $(0,6,0)$. Then $B$ buys
good 2 from $C$, producing $(0,0,6)$. Sales total 12 using a stock of 6.
Reversing the order is infeasible because $B$ initially has no money.
If $B$ hoards the proceeds, the second sale need not occur. Equal
aggregate stocks and quotes therefore permit different realized paths.
This establishes a possibility of interrupted circulation, not a theorem
that interruption or crisis is necessary.

\section{Bridge to the existing mathematical-finance research}
\subsection{What this installment provides}
The preceding models provide distinct objects for physical inventory,
socially normal labor requirements, monetary expression, observed
prices, and realized payments. Later financial contracts must specify
which cash flows they claim, who owes them, and when they are payable.
An inventory labor coefficient alone does not determine a traded claim's
price or an equivalent martingale measure.

The existing draft's interpretation of Girsanov as economic redistribution
requires particular care. A change of probability measure reweights
events for a mathematical representation. It does not itself transfer
goods, money, or surplus-value between owners. Redistribution needs
cash-flow and ownership equations, such as \eqref{eq:trade}, extended
to production and claims.

\subsection{Finite-horizon measure changes: an explicit check}
The existing draft also links asymptotic divergence of a Girsanov kernel
to failure of finite-horizon equivalence. The implication needs an
independent proof and is false in general. Let $W$ be Brownian motion
and take the deterministic kernel $\theta_t=t$. For each finite $T$,
\[
 \E\exp\left(\tfrac12\int_0^T\theta_t^2\dd t\right)
       =\exp(T^3/6)<\infty.
\]
Thus Novikov's sufficient criterion holds on every finite horizon, and
\[
 Z_T=\exp\left(-\int_0^Tt\dd W_t-\tfrac12\int_0^Tt^2\dd t\right)
\]
is strictly positive with $\E Z_T=1$. It defines an equivalent measure
on $\F_T$ even though $\theta_t\to\infty$. Infinite-horizon equivalence
is a different question. Moreover, failure of a \emph{sufficient}
criterion is not proof that no equivalent measure exists. Stochastic
integration and change of measure are treated in \cite{ks}.

Falling profits, default, and liquidity restrictions may be important
economic mechanisms, but their connection to no-arbitrage must be
established for a specified market, admissible strategies, gain
processes, and time horizon. A defaultable payoff is not by itself a
contradiction of arbitrage-free pricing. The counterexample here repairs
the logical bridge for future work; it is not a review of all claims in
the earlier manuscript.

\section{Open closures and the next installments}
The present construction is a baseline, with several open problems.
First, one must choose an interpretation of socially necessary labor:
production norms, market validation, and their interaction need not
produce the same $v$. A validation factor multiplying unsold output
would be an additional theory; neither ``unsold implies zero value''
nor ``all private hours are socially valid'' has been assumed here.

Second, normal techniques and the weights $\alpha_k$ should eventually
be linked to data or institutional mechanisms. Equation
\eqref{eq:reduction} specifies a reduction once weights are given; it
does not determine them. Third, a money-selection model must explain
acceptance rather than obtain money solely by selecting a graph root.
Fourth, price wedges and sale intensities need mechanisms of demand,
access, and bargaining. Those mechanisms determine whether interruptions
are transient, recurrent, or systemic.

To finish the more detailed treatment of Chapter 3, add explicit dated
debts $D^{ac}(T)$, separate commodity delivery from payment, and specify
legally admissible netting and settlement. Hoards can then serve as
payment reserves. For world money, introduce settlement domains,
exchange rates, and international transfers. Converting the accepted
asset to fiat or deposit money requires a balance-sheet issuer and
redemption or settlement rules. These are extensions, not results
already obtained by the closed cash-sale model.

The broader program follows the order of the text:
\begin{enumerate}
\item Volume I, Part II: distinguish $C$--$M$--$C$ from $M$--$C$--$M'$;
formalize labor-power as a special commodity and state the conditions
under which exchange alone cannot explain an aggregate increment.
\item Volume I, Parts III--V: model the labor process, constant and
variable capital, the working day, and absolute and relative surplus-value.
\item Volume I, Parts VI--VIII: wages, reproduction, accumulation, and
the historical conditions of capitalist ownership. Historical arguments
need evidence alongside formal models.
\item Volume II: turnover, inventories, circuit interruption, and
reproduction across departments; distinguish stock and flow dimensions.
\item Volume III: profit formation, competition, commercial and
interest-bearing capital, rent, and financial claims. Reassess the
existing finance results against the foundations established earlier.
\end{enumerate}
The next focused task is the social determination of $\ell$ in Chapter 1,
§§1--2. Comparing a production-normalization closure with an explicit
market-validation closure will expose where their predictions differ
before a more elaborate stochastic system is chosen.

\begin{thebibliography}{9}
\bibitem{marx1}
Karl Marx. \emph{Capital}, Volume I, Part I, Chapter 1,
``Commodities,'' §§1--4. First English edition, 1887; translated by
Samuel Moore and Edward Aveling, edited by Friedrich Engels.
\url{https://www.marxists.org/archive/marx/works/1867-c1/ch01.htm}.

\bibitem{marx2}
Karl Marx. \emph{Capital}, Volume I, Part I, Chapter 2,
``Exchange.'' Same translation.
\url{https://www.marxists.org/archive/marx/works/1867-c1/ch02.htm}.

\bibitem{marx3}
Karl Marx. \emph{Capital}, Volume I, Part I, Chapter 3,
``Money, or the Circulation of Commodities,'' §§1--3. Same translation.
\url{https://www.marxists.org/archive/marx/works/1867-c1/ch03.htm}.

\bibitem{ks}
Ioannis Karatzas and Steven E. Shreve.
\emph{Brownian Motion and Stochastic Calculus}, 2nd edition.
Graduate Texts in Mathematics 113, Springer, 1991.
Chapters 3 and 5. A copy is available in the project's probability library.

\bibitem{prior}
F\'elix Alfonso Camacho Moya. \emph{Stochastic Marxian Macroeconomics:
A Working Draft on Automation, Fictitious Capital, and Networked Crises}.
Undated project manuscript, in the website's \texttt{resources} directory.
Referenced here for continuity, not as independent verification of its
mathematical claims.
\end{thebibliography}
\end{document}
